% !TeX program = xelatex
\documentclass[a4paper,10pt]{article}

% ============================================================================
% សន្លឹកកិច្ចការសិស្ស ស-ខ្មៅ (Student Worksheet - Grade 8 2.1 EXP1)
% បរិបទ៖ គ្មានស្លាយ LCD និងគ្មានសម្ភារៈកាត់បិទ — សិស្សបំពេញផ្ទាល់លើសន្លឹកកិច្ចការ ២ ទំព័រ A4 (១ សន្លឹកសងខាង)
% ============================================================================
\usepackage[a4paper, left=0.85cm, right=0.85cm, top=0.70cm, bottom=0.95cm, footskip=0.48cm]{geometry}
\usepackage{amsmath, amssymb}
\usepackage{unicode-math}
\usepackage{fontspec}
\usepackage{array, colortbl}
\usepackage[table]{xcolor}
\usepackage{enumitem}
\usepackage{pifont}
\usepackage{fancyhdr}
\usepackage{lastpage}
\usepackage{tikz}
\usepackage[most]{tcolorbox}
\usetikzlibrary{calc, arrows.meta, decorations.pathreplacing, positioning, babel}

\makeatletter
\let\@expl@@@mark@update@singlecol@structures@@\relax
\makeatother

% ============================================================================
% ការកំណត់ពុម្ពអក្សរផ្លូវការ (Cross-Platform Fonts)
% ============================================================================
\providecommand{\NGSFindFonts}{%
  \IfFileExists{C:/Windows/Fonts/cambria.ttc}{%
    \def\NGSSysPath{C:/Windows/Fonts/}%
    \def\NGSUserPath{C:/Users/KIT/AppData/Local/Microsoft/Windows/Fonts/}%
  }{%
    \IfFileExists{./fonts/KhmerOSsiemreap.ttf}{\def\NGSSysPath{./fonts/}\def\NGSUserPath{./fonts/}}{%
    \IfFileExists{../fonts/KhmerOSsiemreap.ttf}{\def\NGSSysPath{../fonts/}\def\NGSUserPath{../fonts/}}{%
    \IfFileExists{../../fonts/KhmerOSsiemreap.ttf}{\def\NGSSysPath{../../fonts/}\def\NGSUserPath{../../fonts/}}{%
    \IfFileExists{../../../fonts/KhmerOSsiemreap.ttf}{\def\NGSSysPath{../../../fonts/}\def\NGSUserPath{../../../fonts/}}{%
    \IfFileExists{../../../../fonts/KhmerOSsiemreap.ttf}{\def\NGSSysPath{../../../../fonts/}\def\NGSUserPath{../../../../fonts/}}{%
    \IfFileExists{../../../../../fonts/KhmerOSsiemreap.ttf}{\def\NGSSysPath{../../../../../fonts/}\def\NGSUserPath{../../../../../fonts/}}{%
    \IfFileExists{../../../../../../fonts/KhmerOSsiemreap.ttf}{\def\NGSSysPath{../../../../../../fonts/}\def\NGSUserPath{../../../../../../fonts/}}{%
      \def\NGSSysPath{../../../../../../../fonts/}\def\NGSUserPath{../../../../../../../fonts/}%
    }}}}}}}%
  }%
}
\NGSFindFonts

\setmathfont[Path={\NGSSysPath}, FontIndex=1, Scale=1.02]{cambria.ttc}
\defaultfontfeatures{Ligatures=TeX}

\setmainfont[
  Path={\NGSUserPath}, Extension=.ttf, Script=Khmer,
  BoldFont={KhmerOSsiemreap}, BoldFeatures={FakeBold=1.8},
  ItalicFont={KhmerOSsiemreap}, ItalicFeatures={FakeSlant=0.22}
]{KhmerOSsiemreap}

\newfontfamily\muol[Path={\NGSUserPath}, Extension=.ttf, Script=Khmer]{KhmerOSmuollight}
\newfontfamily\TimesEn[Path={\NGSSysPath}, Extension=.ttf, BoldFont={timesbd}, ItalicFont={timesi}]{times}

\XeTeXlinebreaklocale "km"
\XeTeXlinebreakskip = 0pt plus 1pt minus 0.1pt

\definecolor{panelgray}{gray}{0.93}
\definecolor{headgray}{gray}{0.85}

\pagestyle{fancy}
\fancyhf{}
\renewcommand{\headrulewidth}{0pt}
\fancyfoot[L]{\fontsize{7.5pt}{9.5pt}\selectfont អនុវិទ្យាល័យព្រះក្រសាំង \en{|} គណិតវិទ្យាថ្នាក់ទី៨ \en{|} មេរៀនទី២៖ ២.១. សញ្ញាណស្វ័យគុណ (គ្រូ៖ ហុក ប៊ុនថែន)}
\fancyfoot[R]{\fontsize{7.8pt}{9.5pt}\selectfont ទំព័រ \en{\thepage/\pageref{LastPage}}}

\setlength{\parindent}{0pt}
\setlength{\parskip}{0pt}
\setlength{\tabcolsep}{4pt}

\newcommand{\en}[1]{{\TimesEn #1}}
\newcommand{\enbf}[1]{{\TimesEn\bfseries #1}}
\newcommand{\cnum}[1]{\mbox{\fontsize{10pt}{12pt}\selectfont\ding{\numexpr181+#1\relax}}}
\newcommand{\chkbox}{\tikz[baseline=-0.12ex]\draw[black, line width=0.8pt] (0,0) rectangle (0.24cm,0.24cm);}
\newcommand{\blank}[1][1.5cm]{\makebox[#1]{\dotfill}}
\newcommand{\ansbox}[1][0.85cm]{\,\tikz[baseline=-0.5ex]\draw[black, line width=0.65pt, fill=white] (0,0) rectangle (#1, 0.42cm);\,}
\newcommand{\secbar}[2]{%
  \begin{tcolorbox}[enhanced, sharp corners, colback=headgray, colframe=black, boxrule=0.9pt, boxsep=1pt, top=2pt, bottom=2pt, left=5pt, right=5pt, before skip=4pt, after skip=3pt]
    {\muol\fontsize{9.2pt}{13pt}\selectfont #1}\hfill{\fontsize{8pt}{11pt}\selectfont\textbf{#2}}
  \end{tcolorbox}%
}

\newcolumntype{L}[1]{>{\raggedright\arraybackslash}p{#1}}
\newcolumntype{C}[1]{>{\centering\arraybackslash}p{#1}}

\begin{document}
\fontsize{9.2pt}{13.6pt}\selectfont

% ============================================================================
% ទំព័រទី ១៖ E1 ទស្សន៍ទាយ + E2 រុករក (POE 1 & POE 2) + E3 សំយោគនិយមន័យ និងធនាគារហេតុផល
% ============================================================================
\begin{tcolorbox}[enhanced, sharp corners, colback=white, colframe=black, boxrule=1.0pt, boxsep=2pt, top=3pt, bottom=3pt, left=6pt, right=6pt]
  \begin{center}
    {\muol\fontsize{10.5pt}{15pt}\selectfont អនុវិទ្យាល័យព្រះក្រសាំង \enbf{--} សន្លឹកកិច្ចការសិស្ស \enbf{(}មេរៀនទី២៖ ២\enbf{.}១\enbf{.} សញ្ញាណស្វ័យគុណ\enbf{)}}\\[3pt]
    {\fontsize{8.8pt}{12pt}\selectfont
    ឈ្មោះសិស្ស / ដៃគូ៖ \blank[5.2cm] \hspace{0.25cm}
    ថ្នាក់ទី៨៖ \blank[1.1cm] \hspace{0.25cm}
    ក្រុមទី៖ \blank[1.0cm] \hspace{0.25cm}
    កាលបរិច្ឆេទ៖ ១៦/១១/២០២៦}
  \end{center}
\end{tcolorbox}

% --- E1 + E2 ---
\secbar{ផ្នែកទី១៖ បេសកកម្មទស្សន៍ទាយ \enbf{(E1)} និងការរុករកជាដៃគូ \enbf{(E2 -- POE 1 \& POE 2)}}{រយៈពេល ១៤ នាទី}

\begin{tcolorbox}[enhanced, sharp corners, colback=panelgray, colframe=black, boxrule=0.75pt, boxsep=1.5pt, top=2.5pt, bottom=2.5pt, left=5pt, right=5pt]
  \textbf{បេសកកម្ម \en{E1} (គូសធីក \chkbox\ មុនពិសោធន៍)៖}
  (១) បើបត់ក្រដាសពាក់កណ្ដាល \textbf{៦ ដង} តើបានប៉ុន្មានស្រទាប់? \chkbox\ $12$ ស្រទាប់ \quad \chkbox\ $64$ ស្រទាប់ \quad
  \en{|} \quad (២) តើ $(-2)^6$ និង $-2^6$៖ \chkbox\ ស្មើគ្នា \quad \chkbox\ ផ្ទុយសញ្ញាគ្នា
\end{tcolorbox}

\vspace{2pt}
\textbf{ក. សកម្មភាពទី ១ (\en{POE 1})៖ តារាងបត់ក្រដាសទ្វេដង ធៀបនឹងអន្ទាក់យកគោលគុណនិទស្សន្ត ($2\times n$)}

{\renewcommand{\arraystretch}{1.28}
\begin{tabular}{|C{1.65cm}|L{4.85cm}|C{2.45cm}|C{2.25cm}|C{3.15cm}|C{3.25cm}|}
  \hline
  \rowcolor{headgray}
  \textbf{ចំនួនបត់ $n$} & \multicolumn{1}{c|}{\textbf{ផលគុណកត្តាដូចគ្នា ($\times 2$)}} & \textbf{ទម្រង់ស្វ័យគុណ} & \textbf{ចំនួនស្រទាប់ពិត} & \textbf{អន្ទាក់ខុស ($2\times n$)} & \textbf{ផ្ទៀង ($2^n = 2\times n$?)} \\
  \hline
  $1$ ដង & $2$ ($1$ កត្តា) & $2^1$ & $2$ & $2\times 1 = 2$ & ស្មើចៃដន្យ \\
  \hline
  $2$ ដង & $2\times 2$ ($2$ កត្តា) & $2^2$ & $4$ & $2\times 2 = 4$ & ស្មើចៃដន្យ \\
  \hline
  $3$ ដង & $2\times 2\times 2$ ($3$ កត្តា) & $2^{\blank[0.5cm]}$ & \blank[1.0cm] & $2\times 3 = 6$ & ខុសគ្នា ($8\neq 6$) \\
  \hline
  $4$ ដង & $2\times 2\times 2\times 2$ & $2^{\blank[0.5cm]}$ & \blank[1.0cm] & $2\times 4 = 8$ & ខុសគ្នា ($\blank[0.7cm]\neq 8$) \\
  \hline
  $5$ ដង & \blank[3.8cm] & $2^{\blank[0.5cm]}$ & \blank[1.0cm] & $2\times 5 = 10$ & ខុសគ្នា ($\blank[0.7cm]\neq 10$) \\
  \hline
  $6$ ដង & \blank[3.8cm] & $2^{\blank[0.5cm]}$ & \blank[1.0cm] & $2\times 6 = 12$ & ខុសគ្នា ($\blank[0.7cm]\neq 12$) \\
  \hline
\end{tabular}}

\vspace{3pt}
$\bullet$ \textbf{ពន្យល់ (\en{E})៖} ហេតុអ្វី $a+a+a = \blank[1.1cm]$ ចំណែក $a\times a\times a = \blank[1.1cm]$? ព្រោះស្វ័យគុណជាផល\blank[1.8cm] នៃកត្តាដូចគ្នា!

\vspace{4pt}
\textbf{ខ. សកម្មភាពទី ២ (\en{POE 2})៖ ប្រអប់ចាក់សោគោល \en{(Lock-Box)} និងធ្នូផ្គូផ្គងសញ្ញាដកពីរៗ $\overbrace{(-)(-)}^{\mathbf{+}}$}

{\renewcommand{\arraystretch}{1.34}
\begin{tabular}{|L{2.85cm}|L{3.10cm}|L{7.20cm}|C{4.45cm}|}
  \hline
  \rowcolor{headgray}
  \textbf{កន្សោម} & \textbf{អ្វីជាគោលពិតប្រាកដ?} & \textbf{ពន្លាតកត្តា និងគូសធ្នូផ្គូផ្គងសញ្ញាដកពីរៗ} & \textbf{លទ្ធផល និងសញ្ញា} \\
  \hline
  $\big[-2\big]^6$ & គោលគឺ $(-2)$ (ជាប់ដក) & $\overbrace{(-2)(-2)}^{+4}\times\overbrace{(-2)(-2)}^{\blank[0.8cm]}\times\overbrace{(-2)(-2)}^{\blank[0.8cm]}$ & $=\blank[1.2cm]$ (សញ្ញា $\blank[0.6cm]$) \\
  \hline
  $-2^6$ & គោលគឺ $2$ (ដកនៅក្រៅ!) & $-\big(2\times 2\times 2\times 2\times 2\times 2\big)$ & $=\blank[1.2cm]$ (សញ្ញា $\blank[0.6cm]$) \\
  \hline
  $\big[-3\big]^3$ & គោលគឺ \blank[1.3cm] & $\overbrace{(-3)(-3)}^{\blank[0.8cm]}\times(-3)$ (សល់ដក $1$ កត្តា) & $=\blank[1.2cm]$ (សញ្ញា $\blank[0.6cm]$) \\
  \hline
  $\left(\dfrac{2}{3}\right)^3$ & គោលគឺប្រភាគ $\dfrac{2}{3}$ ទាំងមូល & $\dfrac{2}{3}\times\dfrac{2}{3}\times\dfrac{2}{3} = \dfrac{2\times 2\times 2}{3\times 3\times 3}$ & $=\dfrac{\blank[0.8cm]}{\blank[0.8cm]} \neq \dfrac{8}{3}$ \\
  \hline
\end{tabular}}

% --- E3 ---
\secbar{ផ្នែកទី២៖ សំយោគនិយមន័យ \enbf{(E3)} និងធនាគារហេតុផល \enbf{B\&W} \cnum{1}\enbf{--}\cnum{5}}{សៀវភៅពុម្ព ទំ.១៣–១៤}

\begin{tabular}{|L{9.20cm}|L{9.20cm}|}
  \hline
  \textbf{១. សមាសភាគ និងនិយមន័យស្វ័យគុណ៖}\par\vspace{2pt}
  $\bullet$ ក្នុង $10\,000 = 10^4$៖ ចំនួន $10$ ហៅថា \blank[1.8cm], ចំនួន $4$ ហៅថា \blank[2.0cm], និង $10^4$ ហៅថា \blank[1.9cm] ។\par\vspace{2pt}
  $\bullet$ \textbf{របៀបអាន៖} $5^2$ អានថា $5$ \blank[1.3cm] ; $9^3$ អានថា $9$ \blank[1.3cm] ; $8^4$ អានថា $8$ ស្វ័យគុណ $\blank[0.6cm]$ (កែពុម្ពទំ.១៣ ពី $8^2$ ជា $8^4$) ; សន្មត $a^1 = \blank[0.6cm]$ ។\par\vspace{2pt}
  $\bullet$ \textbf{និយមន័យ (ទំ.១៤)៖} បើ $a\in\mathbb{Q}$ និង $n$ ជាចំនួនគត់\textbf{វិជ្ជមាន} ($n\ge 1$)៖\par\vspace{2pt}
  {\centering $\underbrace{a\times a\times a\times\dots\times a}_{n\text{ កត្តា}} = \ansbox[1.0cm]$\par}\vspace{2pt}
  $\bullet$ \textbf{សម្គាល់៖} $(-1)^n = \blank[0.7cm]$ បើ $n$ គូ \; ; \; $(-1)^n = \blank[0.8cm]$ បើ $n$ សេស ។
  &
  \textbf{២. ធនាគារហេតុផល \en{B\&W} (\cnum{1}--\cnum{5}) សម្រាប់សម្អាង៖}\par\vspace{2pt}
  \cnum{1} \textbf{និយមន័យ៖} $\underbrace{a\times\dots\times a}_{n\text{ កត្តា}} = a^n$ ($n\ge 1$, និង $a^1=a$) ។\par\vspace{2pt}
  \cnum{2} \textbf{របាំងប្រឆាំងអន្ទាក់៖} $a^n \neq a\times n$ និង $a+a+\dots+a = na \neq a^n$ ។\par\vspace{2pt}
  \cnum{3} \textbf{ប្រអប់ចាក់សោ៖} $(-a)^n \neq -a^n$ ; ហើយ $(-a^n) = -(a^n)$ (ដកនៅក្រៅស្វ័យគុណ ដូចជា $(-3^2) = -9$) ។\par\vspace{2pt}
  \cnum{4} \textbf{គោលអវិជ្ជមានក្នុងវង់ក្រចក៖} $n$ គូ $\to (+)$ ; $n$ សេស $\to (-)$ ។\par\vspace{2pt}
  \cnum{5} \textbf{ប្រភាគ \& អាទិភាព៖} $\left(\frac{a}{b}\right)^n = \frac{a^n}{b^n} \neq \frac{a^n}{b}$ ; ធ្វើ \textbf{វង់ក្រចក $\to a^n \to \times/\div \to +/-$} ។ \\
  \hline
\end{tabular}

\vspace{5pt}
\textbf{៣. ហ្វឹកហាត់រហ័សបែងចែកគោល និងនិទស្សន្ត (២ នាទី — បំពេញមុនចូលទំព័រ ២)៖}

{\renewcommand{\arraystretch}{1.35}
\begin{tabular}{|C{3.10cm}|C{3.40cm}|C{2.80cm}|L{8.45cm}|}
  \hline
  \rowcolor{headgray}
  \textbf{ស្វ័យគុណ} & \textbf{គោលពិតប្រាកដ} & \textbf{និទស្សន្ត} & \textbf{របៀបអាន និងការប្រុងប្រយ័ត្នសញ្ញា} \\
  \hline
  $(-5)^4$ & \blank[1.8cm] & \blank[1.2cm] & អានថាដក $5$ ទាំងអស់ស្វ័យគុណ $4$ (លទ្ធផលសញ្ញា $\blank[0.7cm]$) \\
  \hline
  $-5^4$ & \blank[1.8cm] & \blank[1.2cm] & ផ្ទុយនៃ $5$ ស្វ័យគុណ $4$ (សញ្ញាដកនៅក្រៅ $\to$ សញ្ញា $\blank[0.7cm]$) \\
  \hline
  $\left(-\dfrac{4}{3}\right)^7$ & \blank[1.8cm] & \blank[1.2cm] & គោលប្រភាគដកក្នុងវង់ក្រចក និទស្សន្តសេស ($\to$ សញ្ញា $\blank[0.7cm]$) \\
  \hline
\end{tabular}}

\newpage

% ============================================================================
% ទំព័រទី ២៖ E4 លំហាត់គំរូតាមជំហានស៊ើបអង្កេត + E5 ប្រតិបត្តិ & វិនិច្ឆ័យ A-D + Exit Check & កិច្ចការផ្ទះ
% ============================================================================
\secbar{ផ្នែកទី៣៖ ពង្រីកចំណេះដឹងតាមជំហានស៊ើបអង្កេត \enbf{(E4 -- Worked Examples)} ទំ\enbf{.}១៤}{រយៈពេល ១០ នាទី}

\textbf{១. លំហាត់គំរូទី ១ (ទំ.១៤)៖} សរសេរក្នុងទម្រង់ស្វ័យគុណ រួចប្រាប់គោល និងនិទស្សន្ត \textit{(ប្រើ \textbf{[ចង្អុលអូស]} រាប់កត្តា ហាមជំនួសលេខឱ្យអក្សរ $a$)}៖\\[2pt]
\begin{tabular}{@{}L{9.4cm}|L{9.3cm}@{}}
  \textbf{ក.} $\underbrace{4\times 4\times 4\times 4\times 4\times 4}_{\blank[0.6cm]\text{ កត្តា}} = \ansbox[0.95cm]$ $\to$ គោលគឺ \blank[0.8cm], និទស្សន្តគឺ \blank[0.8cm] (\cnum{1})
  &
  \textbf{ខ.} $\underbrace{a\times a\times a\times\dots\times a}_{2010\text{ កត្តា}} = \ansbox[1.15cm]$ $\to$ គោលគឺ \blank[0.8cm], និទស្សន្តគឺ \blank[1.0cm] (\cnum{1})
\end{tabular}

\vspace{4pt}
\textbf{២. លំហាត់គំរូទី ២ (ទំ.១៤ + ប្រៀបធៀបការពារអន្ទាក់ ទំ.២៥)៖} គូសប្រអប់ចាក់សោលើគោល ពន្លាតកត្តា និងគណនាតម្លៃ៖

{\renewcommand{\arraystretch}{1.42}
\begin{tabular}{|L{3.10cm}|L{9.45cm}|C{3.15cm}|C{2.10cm}|}
  \hline
  \rowcolor{headgray}
  \textbf{លំហាត់} & \textbf{ជំហានពន្លាតកត្តា និងផ្គូផ្គងសញ្ញា} & \textbf{តម្លៃលទ្ធផល} & \textbf{ហេតុផល} \\
  \hline
  \textbf{ក.} $2^5$ & $= 2\times 2\times 2\times 2\times 2$ \quad \textit{(ប្រយ័ត្ន៖ មិនមែន $2\times 5 = 10$ ទេ!)} & $= \blank[1.5cm]$ & \cnum{1}, \cnum{2} \\
  \hline
  \textbf{ខ.} $\left(\dfrac{2}{3}\right)^3$ & $= \dfrac{2}{3}\times\dfrac{2}{3}\times\dfrac{2}{3} = \dfrac{2\times 2\times 2}{\blank[1.8cm]}$ & $= \dfrac{\blank[0.8cm]}{\blank[0.8cm]}$ & \cnum{1}, \cnum{5} \\
  \hline
  \textbf{គ.} $(-3)^3$ & $= \overbrace{(-3)\times(-3)}^{+9}\times(-3)$ \quad ($3$ កត្តាជាចំនួនសេស) & $= \blank[1.5cm]$ & \cnum{3}, \cnum{4} \\
  \hline
  \textbf{ឃ.} $(-2)^6$ \textbf{ធៀប} $-2^6$ & $(-2)^6 = \overbrace{(-2)(-2)}^{+4}\cdot\overbrace{(-2)(-2)}^{+4}\cdot\overbrace{(-2)(-2)}^{+4}$ \;\textbar\; $-2^6 = -(2^6)$ & $(-2)^6=\blank[0.8cm]$; $-2^6=\blank[0.8cm]$ & \cnum{3}, \cnum{4} \\
  \hline
  \textbf{ង.} $(4^3)(-3^2)$ \scriptsize(ទំ.២៥) & $= (4\times 4\times 4)\times\big[-(3\times 3)\big] = \blank[1.0cm]\times(\blank[0.9cm])$ & $= \blank[1.5cm]$ & \cnum{3}, \cnum{5} \\
  \hline
\end{tabular}}

\vspace{4pt}
$\bullet$ \textbf{សំណួរឆ្លុះបញ្ចាំង \en{(EIA Looking-Back)}៖} ហេតុអ្វី $(-3)^3 = -3^3 = -27$ តែ $(-2)^6 \neq -2^6$? ចម្លើយ៖ \blank[4.8cm]

% --- E5 ---
\secbar{ផ្នែកទី៤៖ វាយតម្លៃលទ្ធផលសិក្សា \enbf{(E5} \enbf{--} ប្រតិបត្តិ ទំ\enbf{.}១៤ \enbf{+} វិនិច្ឆ័យកំហុស \enbf{+} សំបុត្រចេញ\enbf{)}}{រយៈពេល ១០ នាទី}

\textbf{ក. លំហាត់ប្រតិបត្តិ (សៀវភៅពុម្ព ទំ.១៤ — ធ្វើជាដៃគូ រួចប្ដូរសន្លឹកគ្នាផ្ទៀង \en{Pause-Review})៖} \textit{(សម្គាល់៖ សញ្ញា «$\cdot$» ជំនួសសញ្ញាគុណ «$\times$»)}\\[2pt]
\begin{tabular}{|L{9.20cm}|L{9.20cm}|}
  \hline
  \textbf{១. សរសេរកន្សោមខាងក្រោមក្នុងទម្រង់ស្វ័យគុណ៖}\par\vspace{4pt}
  $\bullet$ \textbf{ក.} $x\cdot x\cdot x\cdot x\cdot x\cdot x = \blank[2.0cm]$ \quad (មាន $\blank[0.7cm]$ កត្តា)\par\vspace{6pt}
  $\bullet$ \textbf{ខ.} $\underbrace{y\cdot y\cdot y\cdot y\cdots y}_{2013\text{ កត្តា}} = \blank[2.2cm]$\par\vspace{5pt}
  $\bullet$ \textbf{គ. (សាកល្បង)៖} $(-5)\cdot(-5)\cdot(-5)\cdot(-5) = \blank[2.0cm]$
  &
  \textbf{២. គណនាតម្លៃនៃស្វ័យគុណខាងក្រោម (បង្ហាញជំហានខ្លីៗ)៖}\par\vspace{4pt}
  $\bullet$ \textbf{ក.} $3^4 = \blank[3.6cm] = \blank[1.5cm]$\par\vspace{6pt}
  $\bullet$ \textbf{ខ.} $\left(-\dfrac{1}{2}\right)^6 = \blank[3.2cm] = \blank[1.5cm]$\par\vspace{6pt}
  $\bullet$ \textbf{គ.} $(-2)^7 = \blank[3.4cm] = \blank[1.5cm]$ \\
  \hline
\end{tabular}

\vspace{4pt}
\textbf{ខ. គូសវង់ជុំវិញជម្រើសត្រឹមត្រូវតែ ១ គត់ ក្នុងចំណោម \en{A, B, C, D} រួចសរសេរលេខហេតុផល \cnum{1}--\cnum{5} នៃជម្រើសខុស៖}

{\renewcommand{\arraystretch}{1.32}
\begin{tabular}{|C{4.40cm}|C{4.40cm}|C{4.40cm}|C{4.40cm}|}
  \hline
  \cellcolor{panelgray}\textbf{ជម្រើស \en{A}} & \cellcolor{panelgray}\textbf{ជម្រើស \en{B}} & \cellcolor{panelgray}\textbf{ជម្រើស \en{C}} & \cellcolor{panelgray}\textbf{ជម្រើស \en{D}} \\
  \hline
  $\left(-\dfrac{1}{2}\right)^6 = \dfrac{1}{64}$ & $2^5 = 10$ & $-2^6 = +64$ & $\left(\dfrac{2}{3}\right)^3 = \dfrac{8}{3}$ \\[4pt]
  \chkbox\ ត្រូវ \quad \chkbox\ ខុស (\blank[0.7cm]) & \chkbox\ ត្រូវ \quad \chkbox\ ខុស (\blank[0.7cm]) & \chkbox\ ត្រូវ \quad \chkbox\ ខុស (\blank[0.7cm]) & \chkbox\ ត្រូវ \quad \chkbox\ ខុស (\blank[0.7cm]) \\
  \hline
\end{tabular}}

\vspace{4pt}
\begin{tcolorbox}[enhanced, sharp corners, colback=white, colframe=black, boxrule=0.85pt, boxsep=2pt, top=3pt, bottom=3pt, left=5pt, right=5pt]
  {\muol\fontsize{8.8pt}{12.5pt}\selectfont សំបុត្រចេញរហ័ស \enbf{(Exit Check --} ធ្វើម្នាក់ឯង ៦០ វិនាទី\enbf{)}៖}\\[2pt]
  (១) $(-3)^4 = \blank[1.5cm]$ ចំណែក $-3^4 = \blank[1.5cm]$ \quad \en{|} \quad
  (២) $\left(-\dfrac{2}{5}\right)^3 = \blank[1.8cm]$ \quad \en{|} \quad
  (៣) $2^4 + 1 = \blank[1.4cm] + 1 = \blank[1.3cm]$
\end{tcolorbox}

\vspace{4pt}
\begin{tcolorbox}[enhanced, sharp corners, colback=panelgray, colframe=black, boxrule=0.85pt, boxsep=2pt, top=3pt, bottom=3pt, left=5pt, right=5pt]
  {\muol\fontsize{8.8pt}{12.5pt}\selectfont កិច្ចការផ្ទះ \enbf{(}សៀវភៅពុម្ព ទំព័រ ២៥\enbf{)}៖}\\[1.5pt]
  $\bullet$ \textbf{លំហាត់ \#1 (ក, ខ, គ, ឃ)៖} សរសេរផលគុណជាស្វ័យគុណ រួចប្រាប់គោល និងនិទស្សន្ត ។\\
  $\bullet$ \textbf{លំហាត់ \#3 (ក, ខ, ឃ)៖} គណនា $(4^3)(-3^2)$ \textit{[ប្រយ័ត្ន៖ $(-3^2) = -(3^2) = -9$]}, $\dfrac{4^3-2^2}{3(2^2)}$, និង $5^3 - 2^4 - 3^3$ ។\\
  $\bullet$ \textbf{លំហាត់ \#4៖} ចតុកោណកែងមានបណ្ដោយ $b^3+1$ និងទទឹង $b^2-1$ ។ គណនាបរិមាត្រចតុកោណកែង បើ $b=3$ ។
\end{tcolorbox}

\vspace{4pt}
{\footnotesize $\bullet$ \textbf{ស្វ័យវាយតម្លៃ និងផ្ទៀងផ្ទាត់ជាមួយដៃគូ៖} \chkbox\ ខ្ញុំបែងចែកដាច់ $(-a)^n$ និង $-a^n$ \quad \chkbox\ ខ្ញុំមិនច្រឡំ $a^n$ នឹង $a\times n$ \quad ហត្ថលេខាដៃគូផ្ទៀង៖ \blank[2.8cm]}

\end{document}
